\section{The defect graphs and the two-body form of the components}
\label{sec:qs-classes}

Let $x$ be a generator and $O=\int\omega\,X_1\cdots X_n$ a monomial local functional of degree $|O|$; sums of monomials are treated by linearity. $J=I_{\mathfrak{psl}(4|4)}-I^{CT}_{\mathfrak{psl}(4|4)}[\varepsilon]-\hbar F_1$ with $\hbar F_1$ a one-loop quantisation through arity one along $x$: the equivariant master equation eq c8-qme modulo $\hbar^2$ holds in arity zero and in arity one along $x$. The finite part of the operators is $\psi=0$; its effect is the law (b) of Prop c8-laws (Section \ref{sec:qs-freedom}). The endpoint is $T=\infty$ for $x\in\Sigma\cup\Sigma^*$ (Prop c8-Tsafe) and any $T\in(0,\infty)$ for $x=K_b$.

\begin{proposition}[The defect vertex at arity one]
\label{prop:qs-defect}
By Lemma c8-transfer, $d_T\,O[T]$ is the $\varepsilon\to0$ limit of the flow, with edges $P(\varepsilon,T)$, of the vertex $(d_{CE}+Q)\mathcal J+\tfrac12\{\mathcal J,\mathcal J\}_\varepsilon+\hbar\Delta_\varepsilon\mathcal J$, $\mathcal J=J+\delta\,O(\varepsilon)$. Its part linear in $\delta$, of arity one along $x$ and of order at most $\hbar$, is the sum of:
\begin{itemize}
\item (V1) the classical vertex $\delta\,(sO)^{(1)}(x)=\pm\delta\{S^{(1)}(x),O\}$, and the arity-zero vertex $\delta\,(Q+\{I,-\})O$ joined to a separate vertex $S^{(1)}(x)$;
\item (V2) the special-edge vertices of order $\hbar^0$: $\delta\big(\{S^{(1)}(x),O\}_\varepsilon-\{S^{(1)}(x),O\}\big)$ (special edge between $S^{(1)}(x)$ and $O$), $\delta\big(\{I,O\}_\varepsilon-\{I,O\}\big)$ with a separate $S^{(1)}(x)$ (between $I$ and $O$), $\{I,S^{(1)}(x)\}_\varepsilon-\{I,S^{(1)}(x)\}$ with a separate $O$ (between $I$ and $S^{(1)}(x)$), $\tfrac12(\{I,I\}_\varepsilon-\{I,I\})$ with separate $O$ and $S^{(1)}(x)$ (between two vertices $I$), and the self-contraction $\delta\,\hbar\Delta_\varepsilon O$ through $K_\varepsilon$ with a separate $S^{(1)}(x)$;
\item (V3) the counterterm vertices of order $\hbar$: $-\delta\big(Q\,\mathrm{ct}^{(1)}_x(O)+\{I,\mathrm{ct}^{(1)}_x(O)\}_\varepsilon+\{S^{(1)}(x),\partial_\delta I^{CT}[\varepsilon](O)\}_\varepsilon+\{I^{CT}_{\mathfrak{psl}(4|4)}[\varepsilon]^{(1)}(x)+\hbar F_1^{(1)}(x),O\}_\varepsilon\big)$; the arity-zero vertices $-\delta\big(Q\,\partial_\delta I^{CT}[\varepsilon](O)+\{I,\partial_\delta I^{CT}[\varepsilon](O)\}_\varepsilon+\{I^{CT}[\varepsilon]+\hbar F_1^{(0)},O\}_\varepsilon\big)$ joined to a separate $S^{(1)}(x)$; and the arity-one and arity-zero defect vertices of the action, of order $\hbar$, joined to $O$.
\end{itemize}
The signs are those of the defect vertex; for odd $O$ those of the proof of Lemma c6-bvexact.
\end{proposition}

\begin{proof}
Expand the defect vertex in $\delta$, in arity and in $\hbar$, with $I_{\mathfrak{psl}(4|4)}=I+S^{(1)}+S^{(2)}$ and $O(\varepsilon)=O-\partial_\delta I^{CT}_{\mathfrak{psl}(4|4)}[\varepsilon](O)$, whose arity-one part is $-\mathrm{ct}^{(1)}_x(O)$ (Def c8-renorm). At order $\hbar^0$ the classical master equation $(d_{CE}+Q)I_{\mathfrak{psl}(4|4)}+\tfrac12\{I_{\mathfrak{psl}(4|4)},I_{\mathfrak{psl}(4|4)}\}=0$ (eq c2-mc) removes the $\delta$-independent bracket parts and leaves their special-edge parts; the $\delta$-linear order-$\hbar^0$ part is $(d_{CE}+Q)O+\{I_{\mathfrak{psl}(4|4)},O\}_\varepsilon$, whose arity-one part is (V1) plus the first special edge of (V2) and whose arity-zero part, joined to $S^{(1)}(x)$ elsewhere in the graph, is the second part of (V1) and the second special edge. $\hbar\Delta_\varepsilon$ of $S^{(1)}(x)$ vanishes: its quadratic terms pair a field with the antifield of another field (Table \ref{tab:qs-terms}, Def c7-vdot gives $(f,f^*)$ for even $x$, whose $K_\varepsilon$ self-contraction is a constant), and its cubic terms carry one antifield, $\lambda^*$, and no $\lambda$. The order-$\hbar$ terms are (V3); $\hbar\Delta_\varepsilon$ of the counterterms is of order $\hbar^2$. This is eq o-defect of the odd design with the $\psi$-terms and the vertex $(Q+\{I,-\})O$ included (audit findings G3, G7).
\end{proof}

\begin{definition}[Counterterm cycles along $x$]
\label{def:qs-ct1}
The counterterm cycles of a pair $(X_a,X_b)$ along $x$ are the chains of Def c10-chain with exactly one vertex given by a term of $S^{(1)}(x)$ placed on the cycle (both cycle legs on edges $P$; along $\xi$: $T_1$ on a line or a cubic term with its $\lambda^*$ leg as output) and no special edge. $\mathrm{ct}^{(1)}_x(X_aX_b)$ is $\hbar\,\mathrm{Sing}_\varepsilon$ of their values (Def c4-ict); the arity-one counterterm of $O$ is $\mathrm{ct}^{(1)}_x(O)=\sum_{a<b}\pm\int\omega\cdots\mathrm{ct}^{(1)}_x(X_aX_b)\cdots$ with the Koszul signs of \eqref{eq:qs-decomp}. Its outputs carry the pair degree minus one; along $\Delta$ and $P_b$ it is that of Def c11-graphs.
\end{definition}

\begin{theorem}[Two-body form of the components along every generator]
\label{thm:qs-classes}
Under the hypotheses above:
\begin{itemize}
\item (i) The $\delta$-linear, arity-one, order-$\hbar$ graphs of the flow of the defect vertex fall into the classes
\begin{itemize}
\item (F) the one-loop flow graphs of the vertices (V1);
\item (C2) the defect cycles of the pairs $(X_a,X_b)$ of $O$ (Definition \ref{def:qs-defect}) with trees of $I$ on their outputs and the spectators;
\item (C0) the graphs whose cycle passes through two letters of $O$ and carries the special edge or the $K_\varepsilon$ contraction of (V2) but no vertex $S^{(1)}(x)$, with $S^{(1)}(x)$ on a tree;
\item (Br) the graphs in which the special edge does not lie on the cycle;
\item (Ac) the graphs whose cycle passes through no letter of $O$, joined to $O$ by a tree;
\item (Ct) the trees with one vertex (V3).
\end{itemize}
\item (ii) (F) is the set of one-loop graphs of the arity-one part of $(sO)[T]$ on the left side of eq c8-rep; (C0), with the arity-zero vertices of (Ct) joined to $S^{(1)}(x)$, tends to the trees of $\hbar L_1(O)^{(0)}$ with one vertex $S^{(1)}(x)$, also on the left side; (Br) tends to $0$ together with the special-edge parts of the counterterm brackets (each pair is $e^{-\varepsilon|p|^2}-1$ at a generic momentum times a renormalised subgraph); (Ac), with the action defect vertices of (Ct), tends to $0$ by the master equation in arity zero and in arity one along $x$.
\item (iii) (C2) and the arity-one vertices of (Ct), together with the counterterm trees of the left side, the trees of $-\partial_\delta I^{CT}[\varepsilon](\{S^{(1)}(x),O\})$ and of $-\mathrm{ct}^{(1)}_x(sO)$, have the limit
\begin{equation}
\begin{aligned}
\hbar L_1(O)^{(1)}(x)=\ &\sum_{1\le a<b\le n}(-1)^{|X_b|(|X_{a+1}|+\cdots+|X_{b-1}|)}\int\omega\,X_1\cdots X_{a-1}\,\Upsilon_x(X_a,X_b)\,X_{a+1}\cdots\widehat{X_b}\cdots X_n\\
&-\big\{\hbar F_1^{(1)}(x),O\big\},
\end{aligned}
\label{eq:qs-decomp}
\end{equation}
the two-body kernel $\Upsilon_x$ of Definition \ref{def:qs-defect} (iv) summed over the unordered pairs of letters, and a one-body term, the bracket with the arity-one finite counterterm along $x$, a derivation of the letters.
\item (iv) The coefficients of $\log\varepsilon$ and of $\varepsilon^{-m}$ cancel. Since the defect cycles contain no $\log\varepsilon$ (Theorem \ref{thm:qs-simplex}), the coefficient of $\log\varepsilon$ gives the arity-one logarithmic covariance along $x$:
\begin{equation}
\mathrm{lc}\big(\{S^{(1)}(x),O\}\big)-\big\{S^{(1)}(x),\mathrm{lc}(O)\big\}=s\,\mathrm{lc}^{(1)}_x(O)-\mathrm{lc}^{(1)}_x(sO)+\big\{L^{(1)}(x),O\big\},
\label{eq:qs-logcov}
\end{equation}
with $L^{(1)}(x)$ the coefficient of $\log\varepsilon$ in $I^{CT}_{\mathfrak{psl}(4|4)}[\varepsilon]^{(1)}(x)$; and for $m\ge1$, with the Koszul signs of \eqref{eq:qs-decomp},
\begin{equation}
\begin{aligned}
&\sum_{a<b}\pm\Big([\varepsilon^{-m}]\mathcal D_x(X_a,X_b)-\sum_{j\ge1}\tfrac{(-1)^j}{j!}\Big(\big\{I,\mathrm{ct}^{(1)}_{x,m+j}(X_aX_b)\big\}_{H^j}+\big\{S^{(1)}(x),\mathrm{ct}_{m+j}(X_aX_b)\big\}_{H^j}\Big)\Big)\\
&\qquad=s\,\mathrm{ct}^{(1)}_{x,m}(O)+\{S^{(1)}(x),\mathrm{ct}_m(O)\}-\mathrm{ct}_m(\{S^{(1)}(x),O\})-\mathrm{ct}^{(1)}_{x,m}(sO),
\end{aligned}
\label{eq:qs-pole}
\end{equation}
the pole identity. $I^{CT}_{\mathfrak{psl}(4|4)}[\varepsilon]^{(1)}(x)$ has no term $\varepsilon^{-m}$ along $\Sigma$, $\Sigma^*$, $\Delta$, $P_b$ and $K_b$.
\item (v) \eqref{eq:qs-decomp} holds for every local functional, gauge-variant ones included; the operator $sO$ enters only through the counterterm identity \eqref{eq:qs-pole} and the arity-zero representative. Along $K_b$ it holds at every $T\in(0,\infty)$ with the same right side.
\end{itemize}
\end{theorem}

\begin{proof}
(i) At order $\hbar$ a graph of the flow has either one cycle and vertices of order $\hbar^0$, or no cycle and one vertex of order $\hbar$ (Prop c3-graphs). The arity-one, $\delta$-linear part places $x$ on exactly one vertex and $\delta$ on exactly one vertex. A graph with a cycle contains exactly one vertex of (V1) or one special edge (or $K_\varepsilon$ contraction) of (V2), the other vertices being $I$, $S^{(1)}(x)$ (once) and $O$ (once). With a vertex of (V1) it is in (F). With a special edge of (V2): if the cycle passes through $O$, it uses two letters of $O$ (Thm c10-decomp (i)); if $S^{(1)}(x)$ and the special edge both lie on it, the graph is in (C2); if the special edge lies on it and $S^{(1)}(x)$ does not, (C0); if the special edge is not on the cycle, (Br); if the cycle avoids $O$, (Ac) or (Br). Trees with one vertex of order $\hbar$ are (Ct).

(ii) The left side of eq c8-rep at arity one along $x$ is $(L(O))[T]^{(1)}(x)=\big(L(O)^{(1)}(x)\big)[T]+\big(L(O)^{(0)}\big)[T]^{(1)}(x)$, the second term being the flow of the arity-zero value with one vertex $S^{(1)}(x)$ (Thm c8-L (iv), the cochain $M[T]$ of Def c8-renorm being $C^\bullet$-linear), and $L(O)=sO+\hbar L_1(O)$. The one-loop part of $(sO)[T]^{(1)}(x)$ consists of the one-loop graphs of the vertex $(sO)^{(1)}(x)=\pm\{S^{(1)}(x),O\}$ and of the arity-zero vertex $(sO)^{(0)}=(Q+\{I,-\})O$ with one $S^{(1)}(x)$, renormalised by their counterterms: these are the graphs (F), and their counterterm trees are the two trees named in (iii) together with the arity-zero counterterms of $sO$ joined to $S^{(1)}(x)$. For (C0): the cycle with the special edge through two letters, without $S^{(1)}(x)$, is an arity-zero defect cycle; at arity zero the same analysis (Thm c10-decomp, eq c10-finitesp, Prop c10-arity1) shows that these cycles with the arity-zero counterterm vertices of (Ct) tend to the tree flow of $\hbar L_1(O)^{(0)}$; the trees carrying $S^{(1)}(x)$ are attached at generic momenta and converge to the corresponding trees with edges $P(0,T)$, the pole parts cancelling as functions of $\varepsilon$ against the same trees on the counterterm vertices. Hence (C0) and these vertices give $(\hbar L_1(O)^{(0)})[T]^{(1)}(x)$, which equals the corresponding part of the left side. (Br): Thm c8-obstruction (ii) and Prop c4-defect1 (the class I-b of DESIGN\_STEP2 Part C). (Ac): the cycle and its counterterm vertices form the flow of the defect of the action in arity zero (with $S^{(1)}(x)$ on a tree) or in arity one along $x$; by Def c8-ob and the master equation in these arities its limit is $0$, and the tree joining it to $O$ converges.

(iii) What remains on both sides is: on the right, (C2) and the arity-one vertices of (Ct); on the left, $\big(\hbar L_1(O)^{(1)}(x)\big)[T]$ at tree level and the two counterterm trees. The vertex $-\{\hbar F_1^{(1)}(x),O\}_\varepsilon$ gives $-\{\hbar F_1^{(1)}(x),O\}$, its special-edge part tending to $0$ since $F_1^{(1)}(x)$ has no poles. The tree vertex $-\{I^{CT}_{\mathfrak{psl}(4|4)}[\varepsilon]^{(1)}(x),O\}_\varepsilon$ is $\log\varepsilon$ times a local functional, by the last statement of (iv), and its special-edge part tends to $0$. The vertices $-\{I,\mathrm{ct}^{(1)}_x(O)\}_\varepsilon$ and $-\{S^{(1)}(x),\partial_\delta I^{CT}[\varepsilon](O)\}_\varepsilon$ split into the brackets without special edge, which are poles and logarithms times local functionals, and their special-edge parts: with $\{f,g\}_\varepsilon-\{f,g\}=\sum_{j\ge1}\frac{(-\varepsilon)^j}{j!}\{f,g\}_{H^j}$, the part of order $\varepsilon^0$ of the special-edge part of $-\varepsilon^{-m}\{f,g_m\}_\varepsilon$ is $-\frac{(-1)^m}{m!}\{f,g_m\}_{H^m}$, which is the finite special vertex \eqref{eq:qs-finitesp}; the terms with $j<m$ are poles and those with $j>m$ tend to $0$; $\log\varepsilon$ times a special edge tends to $0$. Each defect cycle has the expansion of Theorem \ref{thm:qs-simplex}: poles times local functionals, plus $[\varepsilon^0]\mathcal D_x$, a local functional, plus a remainder tending to $0$; the trees on its outputs and the trees on the pole vertices are the same trees with edges $P(\varepsilon,T)$, so after the cancellation of the poles (iv) their difference tends to the trees with edges $P(0,T)$ of the $\varepsilon^0$ coefficients. The limit is therefore the tree flow of the right side of \eqref{eq:qs-decomp}, and Thm c8-L (i) (injectivity of the tree flow) identifies the local functionals. The decomposition over the pairs with the Koszul signs is that of Thm c10-decomp (iii): a cycle passes through $O$ along exactly two letters, the other letters being spectators.

(iv) The left side has a limit, so the local coefficients of $\log\varepsilon$ and $\varepsilon^{-m}$ of the sum vanish. Collecting them: the poles of (C2), the poles of the special-edge parts of the (Ct) brackets shifted by $j$, the brackets without special edge $-s\,\mathrm{ct}^{(1)}_x(O)-\{S^{(1)}(x),\partial_\delta I^{CT}[\varepsilon](O)\}-\{I^{CT}_{\mathfrak{psl}(4|4)}[\varepsilon]^{(1)}(x),O\}$, and $+\partial_\delta I^{CT}[\varepsilon](\{S^{(1)}(x),O\})+\mathrm{ct}^{(1)}_x(sO)$ moved from the left side; this gives \eqref{eq:qs-logcov} and \eqref{eq:qs-pole}. The absence of poles in $I^{CT}_{\mathfrak{psl}(4|4)}[\varepsilon]^{(1)}(x)$: its coefficient of $\varepsilon^{-m}$ is a colour-singlet local functional $\int\sum_px^{\mu_1}\cdots x^{\mu_p}f_{\mu_1\cdots\mu_p}$ of degree $-1$ and total Weyl weight $4-\tfrac12\lambda_x-2m$, the position at weight $-1$ and $p$ at most the degree of $x$ in the position (Prop c10-bound1); a colour singlet of degree $-1$ contains an antifield of degree $-1$ and weight at least $\tfrac52$ and a second adjoint field of weight at least $1$, or $c^*$ (weight $4$) and $c$, so each density $f$ has weight at least $\tfrac72$ and the total weight is at least $\tfrac72-p$. For $m\ge1$ the total weight is at most $\tfrac52$ along $\Sigma$ ($p=0$), $\tfrac32$ along $\Sigma^*$ ($p\le1$), $2$ along $\Delta$ ($p\le1$), $3$ along $P_b$ ($p=0$) and $1$ along $K_b$ ($p\le2$), below $\tfrac72$, $\tfrac52$, $\tfrac52$, $\tfrac72$, $\tfrac32$ respectively (along $\Delta$ this is Thm c9-quadratic (i) and Thm c11-local (ii)).

(v) The argument of (i)--(iv) uses $sO$ only through (F), which cancels between the two sides, and through the counterterm trees of the left side, which enter (iv) only. Along $K_b$, every step holds at fixed finite $T$: the edges $P(\varepsilon,T)$ are bounded uniformly in $T$ and decay as $e^{-|y|^2/4T}$ (Def c8-PT), the graphs of (Br) and (Ac) carry at fixed $\varepsilon$ the term $\tfrac3{8\pi^2}\hbar\log T$ times a tree with the factor $e^{-\varepsilon|p|^2}-1$ (Remark after Thm c8-K), which tends to $0$ at fixed $T$, and $\Upsilon_{K_b}$ is independent of $T$ by Theorem \ref{thm:qs-simplex} (iv); Thm c8-K gives the independence of the local representative.
\end{proof}

\begin{corollary}[Existence of the local representatives through arity one]
\label{cor:qs-exist}
If the one-loop equivariant master equation holds in arity zero and in arity one along $x$, then for every local functional $O$ the local representative of $d_TO[T]$ exists in arity one along $x$ and is given by \eqref{eq:qs-decomp} together with the arity-zero representative $L(O)^{(0)}$; along $\Sigma$, $\Sigma^*$, $\Delta$, $P_b$ and the rotations it is computed at $T=\infty$, along $K_b$ at any finite $T$, with the same result.
\end{corollary}

\begin{proof}
Theorem \ref{thm:qs-classes} exhibits $d_TO[T]$ in arity one along $x$ as the tree flow of the local functional \eqref{eq:qs-decomp}, plus the arity-one flow of $L(O)^{(0)}$; the latter exists by the same argument at arity zero.
\end{proof}

\begin{remark}[Hypothesis (iii) of Def c10-hyp]
\label{rem:qs-hyp}
Corollary \ref{cor:qs-exist} removes the hypothesis that the local representative exists, used in Def c10-hyp (iii), Prop c10-arity1 and Prop c11-oddbound: along every generator it follows from the locality of the $\varepsilon\to0$ limit of each defect cycle, which is Theorem \ref{thm:qs-simplex}, and from the master equation through arity one along $x$. The hypotheses that remain are the master equation, whose arity-one part along $x$ is the equation $sF_1^{(1)}(x)=\mathrm{Ob}^{(1)}(x)$ for the finite counterterm, and, for $T=\infty$, the convergence of Prop c8-Tsafe.
\end{remark}

\begin{corollary}[Dilatation and translations]
\label{cor:qs-DP}
On every local functional $O$ with $sO=0$,
\[
\sum_{a<b}\pm\Upsilon_\Delta(X_a,X_b)=\sum_{a<b}\pm\big(4\,\mathrm{lc}(X_aX_b)+C^{(2)}_\Delta(X_a,X_b)\big),\qquad \sum_{a<b}\pm\Upsilon_{P_b}(X_a,X_b)=\sum_{a<b}\pm C^{(2)}_{P_b}(X_a,X_b),
\] with the composite pair kernels $C^{(2)}_x$ of Def c11-sp (Prop c11-routes, Cor c11-D1local); the one-body terms of the two forms agree.
\end{corollary}

\begin{proof}
Both forms compute the unique local representative (Thm c8-L (i)); for $x=\Delta$ this is Prop c11-routes, whose route of the special edges is \eqref{eq:qs-decomp}; for $x=P_b$ the same argument applies, with Cor c11-D1local and the one-body term $-\{\hbar F_1^{(1)}(P_b),-\}$ in both forms. Pair by pair the two kernels may differ by terms whose sum over the pairs of every cocycle vanishes; \eqref{eq:qs-decomp} holds also for $sO\neq0$ (Theorem \ref{thm:qs-classes} (v)), while the composite form acquires the renormalised composite graphs of $sO$ (Thm c11-local (vi)).
\end{proof}

\begin{remark}[The placements along $\Sigma$, $\Sigma^*$ and $K_b$]
\label{rem:qs-placements}
In the classes (A)--(D) of the odd design (thm o-classes), (C2) consists of: the special edge between $S^{(1)}(x)$ and a letter (A), with $S^{(1)}(x)$ placed at the special edge; between $I$ and $S^{(1)}(x)$ (B), the same placement at an internal vertex; between $I$ and a letter (C), with $S^{(1)}(x)$ on a line, at no special edge, or cubic; between two vertices $I$ (D), likewise. Along $\xi\in S_+\otimes W$ the placements are: $T_1$ on a line (the $\lambda_-$--$c$ composite line of Corollary \ref{cor:qs-edges} (ii)); $T_1,\dots,T_4$ at the special edge (the antifield leg $A^*|_\perp$ of $T_2$ and $\lambda^*$ of $T_3$, $T_4$ only there, since $P$ has no component on them); $T_5$ cubic, with the leg $\lambda^*_\pm$ as an output or on the special edge. Along $\eta$ the same with the Killing spinor $\rho(x)\eta$ and the term $s_0$ at the special edge. Along $K_b$: $(\AV,A^*|_V)$ on a line (the $\AV$--$c$ composite line, Lemma c11-lines (iii)), and each of the five quadratic terms at the special edge. Every graph of (C2) has $k_I\ge0$ vertices of $I$; $k_I=0$ occurs for the placements at the special edge only.
\end{remark}

\begin{remark}[Colour and Koszul signs]
\label{rem:qs-colour}
Lemma c10-colour applies to the defect cycles with the vertex $S^{(1)}(x)$ read as follows: a quadratic term is a pairing $(-,-)$ of two adjoint fields and acts on the colour of the cycle as the identity; a cubic term $T_5$ has the colour form $([Z_1,Z_2],Z_3)$ and acts as $\mathrm{ad}_Y$ with $Y$ its output, like a cubic vertex of $I$; the special edge carries $(-,-)^{-1}$ like an edge of $P$. Hence the colour factor of a defect cycle in the colour slots of the two letters is $(1\otimes\mathrm{ad}_{Y_r}\cdots\mathrm{ad}_{Y_1})(-,-)^{-1}$ with the outputs $Y_1,\dots,Y_r$ in the order of the chain from $X_a$, the output of $T_5$ included, and the planar reading of Lemma c10-colour (ii)--(iv) applies unchanged: a term $a\frac{\hbar}{16\pi^2}$ times this element contributes $(-1)^r\tfrac a2$ units to an adjacent pair of a single trace. The Koszul signs of \eqref{eq:qs-decomp} are those of Thm c10-decomp (iii); for odd $x$ the component $\hbar L_1(O)^{(1)}(x)$ has the parity of $O$ shifted by the parity of the evaluation, as in Def c8-components.
\end{remark}
